\documentclass[11pt,letterpaper]{article}
\usepackage[margin=1in]{geometry}
\usepackage{amsmath,amssymb}
\usepackage{booktabs}
\usepackage{setspace}
\usepackage{natbib}
\usepackage[hidelinks]{hyperref}
\usepackage{array}
\newenvironment{adjustwidth}%
  {\list{}{\leftmargin=0.09\textwidth \rightmargin=0.09\textwidth \listparindent=1em \itemindent=0pt \parsep=0.5em}\item[]}%
  {\endlist}
\usepackage[protrusion=true,expansion=false]{microtype}          % better justification, far fewer overfull lines
\widowpenalty=10000             % never strand a single line at a page top
\clubpenalty=10000              % never strand a single line at a page bottom
\displaywidowpenalty=10000
\raggedbottom                   % do not stretch pages to fill; keeps spacing even
\usepackage{longtable}
\usepackage{caption}
\captionsetup[table]{skip=6pt,font=small,labelfont=bf}
\renewcommand{\arraystretch}{1.15}
\usepackage{url}
\usepackage{float}
\setcounter{topnumber}{2}
\setcounter{totalnumber}{3}
\renewcommand{\topfraction}{0.85}
\renewcommand{\textfraction}{0.10}
\renewcommand{\floatpagefraction}{0.75}
\bibpunct{(}{)}{;}{a}{}{,}
\onehalfspacing
\setlength{\parskip}{0.35em}

%% ANONYMOUS=1 builds the double-anonymous version for Management Science.
%% Leave undefined for the SSRN version, which REQUIRES author + affiliation.
\newif\ifanon\anonfalse   %% <-- change to \anontrue for the MS submission
\newcommand{\AUTHORNAME}{\ifanon [omitted for review]\else Dhanya MD\fi}
\newcommand{\AUTHORAFFIL}{\ifanon\else Independent Researcher\fi}
\newcommand{\AUTHOREMAIL}{\ifanon\else dhanya13md@gmail.com\fi}

\title{\textbf{Every Asset Its Own Benchmark:\\[2pt]
Market-Neutral Alpha in Perpetual Futures}}
\author{\AUTHORNAME\\ \small \AUTHORAFFIL\\ \small \texttt{\AUTHOREMAIL}}
\date{\vspace{-6pt}\small \today}

\begin{document}
\maketitle
\vspace{-14pt}

\begin{center}
{\scshape Abstract}
\end{center}
\vspace{-4pt}
\begin{adjustwidth}
\singlespacing\small\noindent
Cross-sectional factor investing assumes assets can meaningfully be ranked
against one another. Equities supply the anchor that makes this reasonable:
shared accounting, industry structure, cash flow. Cryptocurrency perpetual
futures supply none of it. We ask what happens to factor construction without it. Transporting ten factors
from the published equity literature to 112 perpetual contracts over a 363-week
panel, we find that ranking each contract against its own trailing history
dominates ranking it against its peers in nine of eleven cases, raising mean
annualised Sharpe from 1.01 to 1.45; the effect persists with the carry component
removed entirely. Standardising each asset against its own history before ranking
cross-sectionally recovers 41\% of the gap, identifying scale heterogeneity as one
channel. We further show that
preserving factor independence through construction, by forming one book per factor
and combining returns rather than scores, adds 0.665 Sharpe ($t=5.82$) when both
constructions are tuned over an identical grid. Regressed on market, size and
momentum factors the strategy retains an annualised alpha of 27\%, with a market
loading indistinguishable from zero. Through the 2022 decline, in which the
market fell 82\%, it returned $+9\%$ with a maximum drawdown of 8\%. We introduce a factor exploiting aggressor-side information that
centralised venues publish exactly and equity research must infer. Results
survive walk-forward selection across seven rebalance anchors, a held-out
universe test, costs at four times the modelled level, and a deflated Sharpe
correction whose trial count is read from an executed log.

\vspace{6pt}
\noindent\textbf{Keywords:} factor investing, cross-sectional ranking,
cryptocurrency, perpetual futures, portfolio construction, backtest overfitting

\vspace{2pt}
\noindent\textbf{JEL:} G11, G12, G14, C58

\vspace{4pt}
\noindent\textbf{Disclosure.} AI tools were used in this research and for
re-verification of the analysis code and tests. All
results were independently verified by re-execution, and the author is
responsible for all claims and interpretations.

\end{adjustwidth}
\normalsize
\onehalfspacing

\thispagestyle{empty}
\newpage
\setcounter{page}{1}

\section{Introduction}

Every cross-sectional factor model begins with a comparison. We say an asset is
cheap, or small, or heavily traded, and each term carries an implicit completion:
\emph{relative to its peers}. The comparison is defensible in equities because
the market supplies scaffolding that makes it meaningful. Firms report under
common accounting standards, occupy classifiable industries, and are valued
against cash flows that exist. Ranking them compares like with like.

Cryptocurrency perpetual futures dissolve that scaffolding. There are no
earnings, no industry classification, and no accounting identity tying price to
anything. Whether one contract's turnover is high \emph{relative to another's} is
a question the data cannot really answer, because the quantities were never
commensurate. Yet the factor literature, in turning to this market, has largely
imported the cross-sectional apparatus unchanged.

This paper asks what that assumption costs, and what replaces it.

Our design is deliberately conservative. Rather than proposing new signals, we
transport ten factors established in the published equity literature to 112
perpetual contracts across a 363-week panel, retaining their original definitions
and sign conventions. Nothing is estimated. The only quantity we vary is the
frame of comparison: in one specification each contract is ranked against its
peers within the week, as the literature does; in the other, against its own
trailing history, asking not whether an asset is unusual relative to others but
whether it is unusual \emph{for itself}.

The eleven factors are the \emph{sample}, not the contribution. A claim about
construction demonstrated on a single factor would be a curiosity, indistinguishable
from a fortunate choice of signal; demonstrated across eleven factors drawn from
the published literature and spanning price, flow, positioning and trade-size
distribution, it is evidence about the method rather than about any particular
signal. We accordingly report per-factor results throughout and never a single
headline attributable to one of them.

\textbf{First, the comparison frame matters.} Ranking each contract against its
own history outperforms ranking it against peers for nine of the eleven factors
we examine, improving mean risk-adjusted performance by roughly forty per cent.
The advantage is not an artefact of the funding component, which we find to be a
first-order signal in its own right: repeating the comparison with funding
removed entirely leaves the ordering intact.

\textbf{Second, we identify a mechanism rather than only an effect.} An
intermediate specification, standardising each asset against its own history
before ranking cross-sectionally, recovers roughly two-fifths of the advantage,
identifying scale heterogeneity as one channel. The remainder survives
standardisation, and we attribute it to distributional robustness: a rank
transform is invariant to the shape of the distribution it is applied to, whereas
standardisation normalises only the first two moments and implicitly assumes
approximate normality. Crypto factor distributions violate that assumption
severely.

\textbf{Third, construction dominates selection.} Standard practice blends factor
scores into a composite and forms one portfolio from the induced ordering. This
compresses several near-independent signals into a single ranking and discards
their independence before any position is taken. Forming one book per factor and
combining the resulting \emph{returns} instead improves performance substantially
and reliably, favourable in fifteen of seventeen matched comparisons. Both
constructions were tuned over an identical grid; comparing a tuned method against
an untuned baseline is the most common way this literature manufactures a result,
and we take care not to.

\textbf{Fourth, we exhibit a factor whose construction depends on information
this market publishes and equity research must infer.} Work on Chinese A-shares
established that predictive content resides in the \emph{shape} of the per-bar
trade-size distribution rather than its mean. That measure is necessarily
direction-blind, since A-share tick data does not reliably identify the
initiating side. Perpetual futures publish taker-side volume on every bar,
permitting the distribution to be partitioned by initiating side and its
asymmetry measured. The economic motivation is \citet{kyle1985}: an informed
trader must transact in size to profit, and trades off concealment against
execution speed, leaving a detectable footprint rather than none. We do not
overstate the novelty: US equity researchers can
infer direction via the tick rule \citep{leeready1991}, with known error rates.
The distinction is measured versus estimated. The resulting factor ranks third of
the eleven by its marginal contribution to the portfolio.

Throughout we hold ourselves to the significance standard of \citet{harvey2016},
who argue that the proliferation of published factors implies a hurdle of $t>3.0$
rather than the conventional $t>2.0$. We report where we clear it and where we do
not. Appendix~A summarises every statistical test applied in this
paper and its outcome; the tests themselves are developed in Sections~5--7.


Finally, this paper reports an audit of its own claims. During preparation we
searched the research codebase for the executable artifact behind every headline
figure; four had none, including a pair of statistics that proved to be
tail-index values transcribed from a document about an unrelated strategy.
Section~\ref{sec:audit} reports this alongside the infrastructure now preventing
recurrence. We regard the reliability of a reported backtest as an empirical
property of the pipeline producing it, and therefore as something to be tested.

Three elements of this study appear to be without precedent in the published
literature, and we state them as such while noting that we have not conducted an
exhaustive search.

\emph{First}, we are not aware of prior work that uses the \emph{skewness of the
trade-size distribution, measured separately by initiating side and differenced
over time}, as a cross-sectional return predictor. Two adjacent literatures come
close without meeting. \citet{zhou2008} partitions the size distribution by side
in an order-driven equity market, but measures \emph{tail exponents} and finds
them statistically indistinguishable across sides ($\alpha=2.30$ versus $2.36$).
Separately, a large literature prices the skewness of \emph{returns}. Neither
measures the skewness of \emph{trade sizes}, and neither examines how its
between-side asymmetry evolves.

The distinction matters for interpretation as well as for priority. Zhou's null
result says the two sides' size distributions are equally fat-tailed; our result
says their \emph{shapes shift at different times}, and that those shifts predict.
The two are consistent: a distribution can have a stable tail exponent while its
skewness moves.

\emph{Second}, the requirement that the deflation inputs be measured rather than
assumed is understood in principle, the dispersion term cannot be formed without
the full set of trial Sharpes, but it is rarely satisfied in practice, because
the trial set is not usually retained. We report a study in which both inputs are
read from an append-only log written by the code performing each evaluation, and
we report the resulting value under two defensible definitions of the trial set,
including one that does not clear the threshold.

\emph{Third}, we are not aware of a published strategy paper that verifies the
deployed system computes the same positions as the research code. We report an
automated equivalence gate that does so, and we report that it previously
\emph{failed}; research and production agreed perfectly while both called an
incorrect function. The claim that a live system implements a published backtest
is, on our evidence, an empirical proposition rather than a safe assumption.

We also combine two robustness procedures that to our knowledge have not been
applied together: walk-forward configuration selection and rebalance-timing-luck
variation. Applied separately, each leaves a channel open, a walk-forward result
still inherits the arbitrary choice of rebalance weekday, and a timing-luck study
still inherits a fixed configuration. Section~6.1 closes both simultaneously.

The rest of the paper is organised as follows. Section~2 places the study in the
literature. Section~3 describes the data and the two properties of the venue that
make the exercise possible. Section~4 sets out the ranking primitives, portfolio
formation and cost model. Section~5 reports the main results. Section~6 subjects
them to five independent tests. Section~7 estimates capacity.
Section~8 reports the claims that did not survive our own audit, and Section~9
concludes.

\section{Related literature}

\textbf{Cross-sectional factor construction.} The characteristic-sorted portfolio
is the workhorse of empirical asset pricing \citep{famafrench1993}, and the debate
over whether characteristics or covariances drive returns \citep{danieltitman1997}
presumes throughout that assets are commensurate enough to sort. We take that
presumption as the object of study rather than as background.

\textbf{Factor transport across asset classes.} \citet{asness2013} show that
value and momentum earn premia consistently across eight markets and asset
classes, establishing that factor logic travels beyond the equities in which it
was discovered. We transport factors in the same spirit, but ask a question their
design does not: whether the cross-sectional \emph{frame} travels with them.

\textbf{Why we rank rather than forecast.} \citet{gu2020} document that
single-asset return prediction is a low signal-to-noise problem even with
flexible models. Our book never forecasts a return; it ranks contracts and holds
a dollar-neutral spread, so only the ordering must be correct.

\textbf{Time-series versus cross-sectional signals.} \citet{moskowitz2012}
established that an asset's own past return predicts its future return
independently of any cross-sectional ordering. Our ranking primitive is the same
statistical object; the contribution is not the primitive but its systematic
application across eleven factors spanning price, flow, positioning and
trade-size distribution, with a decomposition of why it works.

\textbf{Cryptocurrency asset pricing.} \citet{cong2026} construct a four-factor
cross-sectional model for digital assets, introducing a value factor built from
the ratio of active addresses to market capitalisation alongside market, size and
momentum factors, and document pronounced segmentation across token categories.
Their design is instructive for our purposes in two ways. First, it is
\emph{cross-sectional}, which makes Section~\ref{sec:frame} a direct empirical
response rather than a parallel exercise: we find that the same ranking frame
they employ is materially weaker in this market than ranking each contract
against its own history. Second, and more interestingly, their value factor is an
attempt to \emph{manufacture} the comparability anchor that equities possess
natively, active addresses stand in for the fundamental quantity crypto does not
report. That such a construction is necessary is itself evidence for the problem
we identify. Our response differs in kind: rather than synthesising an anchor, we
abandon cross-asset comparison where it is not warranted. The two approaches are
complementary, and we regard testing their factors under a self-referential frame
as the natural next step.

\textbf{Backtest overfitting.} \citet{bailey2014} formalise deflation of an
observed Sharpe by the expected maximum across trials; \citet{harvey2016} raise
the significance hurdle; \citet{hoffstein2019} document that rebalance-schedule
choice alone generates dispersion often exceeding the effect under study. We
apply all three, and in the first case address the standard weakness that the
trial count is supplied by the researcher.

\section{Data}

\textbf{Universe.} 112 USDT-margined perpetual futures contracts, screened for
continuous trading and finite positive prices, spanning 363 weeks. After factor
warm-up the strategy is evaluable on 308 weekly rebalances and active on 282
(Table~\ref{tab:data}).

\begin{table}[!htbp]\centering
\caption{Sample and universe.}\label{tab:data}
\begin{tabular}{ll}
\toprule
Universe & 112 USDT-margined perpetual futures \\
Panel & 363 weeks \\
Rebalances evaluable / active & 308 / 282 \\
Rebalance frequency & weekly \\
Factors & 11 (10 transported, 1 introduced) \\
Cost model & liquidity-scaled maker fee, 1--5 bp \\
Funding & subtracted at every rebalance \\
\addlinespace[2pt]
\bottomrule
\end{tabular}
\end{table}

\textbf{Sources.} Price and volume derive from hourly bars in a public bulk
archive, resampled weekly. Intraday inputs are reduced from 1-hour and 5-minute
bars to one record per UTC day. Positioning data, open interest, top-trader and
all-account long/short ratios, comes from a keyless exchange metrics endpoint at
daily frequency. No proprietary or paid data is used.

\textbf{Two venue properties are load-bearing.} First, aggressor side is
published: every bar carries taker-buy volume, where equity microstructure
research must infer direction \citep{leeready1991}. Second, funding is observable
and material. Measured on the deployed book over 282 weeks, the funding term
averages $-0.173\%$ per week, or $-8.99\%$ annualised, and is a net \emph{receipt}
in 98\% of weeks: because the construction shorts high-funding contracts and
holds low-funding ones, the book is a carry collector rather than a payer. A
price-only backtest therefore \emph{understates} this strategy by approximately
9\% annually. We account for realised funding at every rebalance rather than
approximating it.

\textbf{A resolution constraint discovered by failure.} Trade-size distribution
factors cannot be computed hourly: a day divides into roughly twelve bars per
side, below the twenty-observation minimum the shape estimator requires. Measured
computability is 0\% at 1-hour and 100\% at 5-minute. The remaining intraday
factors measured \emph{better} hourly. The pipeline therefore runs at two
resolutions simultaneously, a constraint the data imposed, not an optimisation.

\section{Methodology}

\subsection{Factors}
\label{sec:factors}

Every factor is specified a priori from its source, including its sign. Nothing
is estimated. This determines what can be overfit: with no fitted parameters
anywhere in the model, the entire overfitting surface reduces to a handful of
discrete hyperparameters, and Section~6.3 holds all of them out. Each raw factor
is smoothed with a twenty-period rolling mean before ranking.

Table~\ref{tab:factors} gives definitions. Three have direct academic
antecedents: the signed jump measure follows the realised-semivariance
decomposition of \citet{bns2010}, whose separation of upside from downside
variation \citet{pattonsheppard2015} show carries predictive content; the
order-flow imbalance follows \citet{crs2002}; and the trade-size distribution
factors build on the literature discussed in Section~4.4. The remainder are
reproductions of Chinese sell-side research on A-shares. We note plainly that
those are practitioner research reports rather than peer-reviewed publications,
and are not always formally citable; we reproduce the constructions as described
and report their performance without relying on the original claims.

One sign convention is inverted relative to its source. The smart-money factor is
documented as a reversal signal in A-shares; in perpetual futures it is
momentum-consistent. We report this as an empirical finding rather than a fitted
choice, and record it as a limitation: a single inverted sign, chosen with
knowledge of the data, is a channel through which information can leak.

\begin{table}[!htbp]\centering
\caption{Factor definitions. All inputs are point-in-time. $v_n$, $n_n$ and $b_n$
denote volume, trade count and taker-buy volume in intraday bar $n$; $c_n$ its
close. Weekly quantities are smoothed with a twenty-period rolling mean before
ranking.}\label{tab:factors}
\small
\begin{tabular}{p{0.09\textwidth}p{0.40\textwidth}p{0.20\textwidth}p{0.19\textwidth}}
\toprule\addlinespace[2pt]
factor & construction & economic reading & origin \\
\midrule\addlinespace[2pt]
AVOL & $-\log \sum_{12w} \text{volume}$ & abnormal turnover predicts negatively & A-share practitioner \\
Q & VWAP of the top 20\% of volume by $|r|/\sqrt{v}$, over VWAP of all bars & informed accumulation below the average price & A-share practitioner \\
RSJ & $-(\sum_{r>0} r^2 - \sum_{r<0} r^2)/\sum r^2$ & upside-jump-dominated names underperform & \citet{bns2010} \\
OFI & $(2\sum b_n - \sum v_n)/\sum v_n$ & net aggressive buying persists & \citet{crs2002} \\
CPVm & $-\,$mean of daily corr$(c_n, v_n)$ & price and volume moving together signals crowding & A-share practitioner \\
CPVv & $-\,$s.d.\ of daily corr$(c_n, v_n)$ & instability in that relation signals regime change & A-share practitioner \\
WRspread & $-(\text{top-trader L/S} - \text{all-account L/S})$ & large accounts crowded relative to the rest & positioning data \\
TopChg & $-\Delta(\text{top-trader L/S})$ & large accounts \emph{increasing} exposure is a fade & positioning data \\
Quad & $-\operatorname{sign}(r_w)\cdot\Delta\log(\text{open interest})$ & price and open interest rising together is new crowding & CTA practitioner \\
TKU & kurtosis of $\log(v_n/n_n)$ & concentration of the trade-size distribution & A-share practitioner \\
\textbf{TSKD} & $\Delta\big[\mathrm{skew}(\log \text{ticket}_{\mathcal{B}}) - \mathrm{skew}(\log \text{ticket}_{\mathcal{S}})\big]$ & \textbf{shift in which side places the outsized trades} & \textbf{introduced here} \\
\addlinespace[2pt]
\bottomrule
\end{tabular}
\normalsize
\end{table}

\subsection{Directional trade-size asymmetry}
\label{sec:tskd}

Consider one intraday bar containing $n_n$ trades totalling $v_n$ units of
volume. Its \emph{ticket size} is $v_n/n_n$, the typical size of an individual
trade within it. Across a day these form a distribution, and the A-share
literature establishes that its \emph{shape} carries predictive content that its
mean does not.

The limitation of that measure is that a right-skewed distribution is produced
equally by a large buyer and a large seller: it detects that unusually sized
orders occurred, but not who placed them. With taker-buy volume $b_n$ published
per bar, the distribution can be partitioned by initiating side,
\[
\mathcal{B} = \{n : b_n > \tfrac{1}{2}v_n\}, \qquad
\mathcal{S} = \{n : b_n \le \tfrac{1}{2}v_n\},
\]
and its asymmetry measured as
\[
\mathrm{TSKD}_t \;=\; \Delta\Big[\mathrm{skew}\big(\log \mathrm{ticket}_{n\in\mathcal{B}}\big)
\;-\; \mathrm{skew}\big(\log \mathrm{ticket}_{n\in\mathcal{S}}\big)\Big].
\]

A positive value indicates that unusually large trades cluster on the buying
side. The \emph{change} carries the signal where the level does not: what
predicts is not which side is currently large but which side has \emph{just
become} large. The economic motivation is \citet{kyle1985}: an informed trader
must transact in size to profit and trades off concealment against execution
speed, leaving a detectable footprint rather than none.

The measure requires five-minute bars. At hourly resolution a day divides into
roughly twelve bars per side, below the twenty-observation minimum the shape
estimator requires; measured computability is 0\% at one hour and 100\% at five
minutes across 80{,}000 day-records.

\subsection{Ranking primitives}

Let $f_{i,t}$ denote a raw factor observation for asset $i$ at time $t$. The
cross-sectional score ranks within a date,
\[
s^{\mathrm{XS}}_{i,t} = \mathrm{rank}\!\left(f_{i,t} \mid \{f_{j,t}\}_{j=1}^{N}\right),
\]
meaningful only if the $f_{j,t}$ are commensurate. The self-referential score
ranks within an asset,
\[
s^{\mathrm{TS}}_{i,t} = \mathrm{rank}\!\left(f_{i,t} \mid \{f_{i,\tau}\}_{\tau=t-w}^{t-1}\right),
\]
making no cross-asset comparison. Both map to $[-1,+1]$ via
$((\mathrm{rank}-0.5)/n - 0.5)\cdot 2$, which is centred, antisymmetric and
NaN-preserving. We use $w=52$ weeks with a 26-week minimum.

The two degrade differently. As the cross-section becomes more heterogeneous,
$s^{\mathrm{XS}}$ increasingly ranks \emph{asset identity} rather than
\emph{signal}: a contract whose order flow is systematically larger than
another's occupies the same tail week after week irrespective of new information.
$s^{\mathrm{TS}}$ is immune by construction.

\subsection{Portfolio formation}

Each factor independently forms a dollar-neutral quintile book, long the top
20\%, short the bottom 20\%, weights $\pm 1/n$. Scores are never blended. Book
\emph{returns} are combined by inverse trailing volatility,
$w^{(k)}_t \propto 1/\hat\sigma^{(k)}_{t-1}$ normalised to sum to one, with the
estimate lagged so no contemporaneous information enters. Where the risk model is
not estimable the book is held flat rather than traded equally weighted.

\subsection{Carry and turnover}

Each score is penalised by the contract's cross-sectional funding rank, the one
place a cross-sectional comparison is appropriate, since funding is denominated
identically across contracts. Section~\ref{sec:carry} shows this component is far
more consequential than the term ``tilt'' suggests.

A uniform partial adjustment caps per-rebalance turnover. At the modelled maker
fee this \emph{costs} performance monotonically; what it purchases is robustness
to cost misestimation. The deployed configuration therefore uses a cap that is
not Sharpe-maximising, a deliberate divergence we state rather than leave to be
discovered.

\subsection{Return construction}
\[
r_t = \underbrace{w_t' \, \mathrm{fwd}_t}_{\text{price}}
- \underbrace{w_t' \, \mathrm{funding}_{t+1}}_{\text{carry paid}}
- \underbrace{|w_t - w_{t-1}|' \, c}_{\text{transaction cost}}
\]
where $c$ is a liquidity-scaled maker fee from 1 basis point for the most liquid
quintile to 5 for the least, interpolated on cross-sectional dollar-volume rank.
Forward returns are clipped at $+100\%$ so a single listing spike cannot
manufacture a Sharpe.

\subsection{Point-in-time discipline}

Every factor input is leak-tested: a historical score is recomputed with all
future observations deleted and required to reproduce exactly. This test exists
because it previously \emph{failed}. A ranking helper applied along the wrong
array axis ranked each observation against the asset's entire history, future
included, altering historical scores for 99 of 112 contracts (maximum change 1.15). Research and
production code agreed perfectly throughout, because both called the same
incorrect function. All results here postdate the correction and pass at
$0.00\mathrm{e}{+}00$.

\section{Empirical results}

\subsection{The comparison frame}
\label{sec:frame}

Table~\ref{tab:frame} reports eleven single-factor books under both ranking
frames, holding data, universe, costs and construction fixed, both as the
strategy trades it and with carry removed entirely.

As traded, mean annualised Sharpe rises from 1.011 to 1.451; an improvement of
roughly forty-four per cent obtained without changing a single input, only the
set against which each observation is ranked. Nine of eleven factors improve.
With carry removed the improvement is larger in proportional terms and eight of
eleven factors improve. \textbf{The ranking effect is independent of carry}: the
component amplifies the gap but does not create it. No factor changes sign under
either specification, so the appropriate claim is degradation of the
cross-sectional frame, not its failure. One factor reverses the sign of its
ranking preference when carry is removed; at this noise level individual
reversals should not be interpreted, and only the aggregate is meaningful.

\begin{table}[!htbp]\centering
\caption{Ranking frame. Eleven single-factor books, annualised Sharpe, net of
funding and costs. ``Carry removed'' sets the funding weight to zero.}
\label{tab:frame}
\begin{tabular}{lcccc}
\toprule
specification & cross-sectional & self-referential & gap & self wins \\
\midrule\addlinespace[2pt]
as traded (with carry) & 1.011 & \textbf{1.451} & $+0.440$ & 9 of 11 \\
carry removed & 0.383 & \textbf{0.648} & $+0.265$ & 8 of 11 \\
\midrule\addlinespace[2pt]
sign reversals & \multicolumn{4}{c}{\textbf{0 of 11} under either specification} \\
\addlinespace[2pt]
\bottomrule
\end{tabular}
\end{table}

\subsection{Mechanism}

Table~\ref{tab:mech} introduces an intermediate specification: standardise each
asset against its own trailing window, then rank cross-sectionally. If scale
heterogeneity drives the effect this should recover it. It recovers 41\%.

Scale heterogeneity is therefore one channel but not the dominant one. We
attribute the residual to distributional robustness: standardisation normalises
only the first two moments, whereas a rank against an asset's own history is
invariant to the entire distributional shape. Tests separating the fat-tail
component from higher moments are left to future work.

\begin{table}[!htbp]\centering
\caption{Mechanism decomposition.}\label{tab:mech}
\begin{tabular}{lcc}
\toprule
specification & mean annualised Sharpe & gap recovered \\
\midrule\addlinespace[2pt]
cross-sectional, raw & 1.011 & , \\
cross-sectional, per-asset standardised & 1.192 & 41\% \\
self-referential & \textbf{1.451} & 100\% \\
\addlinespace[2pt]
\bottomrule
\end{tabular}
\end{table}

\subsection{Carry is a factor, not an adjustment}
\label{sec:carry}

Table~\ref{tab:carry} reports what removing the funding component costs each
book. The effect is large for every factor: the least affected gains roughly
forty per cent from its inclusion, the most affected more than triples. Traded on its own the component is a respectable strategy in its own right,
though it falls short of the significance threshold we adopt. Carry is therefore
a first-order signal in perpetual futures rather than an accounting adjustment.
Every single-factor figure in this paper includes it, and we say so rather than
presenting those numbers as the ranking effect in isolation.

\begin{table}[!htbp]\centering
\caption{Carry is a factor. Single-factor books with and without the funding
component, matched pairwise so each row uses the largest sample valid for that
comparison.}\label{tab:carry}
\begin{tabular}{lccc}
\toprule
factor & with carry & without carry & carry contributes \\
\midrule\addlinespace[2pt]
TSKD & 1.442 & 0.394 & $+1.048$ \\
TKU & 1.494 & 0.674 & $+0.820$ \\
OFI & 1.241 & 0.480 & $+0.762$ \\
AVOL & 1.222 & 0.663 & $+0.559$ \\
Quad & 1.896 & 1.356 & $+0.540$ \\
\addlinespace[2pt]
\bottomrule
\end{tabular}

\vspace{4pt}
\begin{minipage}{0.80\textwidth}\footnotesize
\emph{Notes.} Traded on its own over the matched sample, the carry component
returns an annualised Sharpe of $1.250$ ($t=2.91$), short of the $t>3.0$
threshold adopted throughout.
\end{minipage}
\end{table}

\subsection{Construction dominates selection}

Table~\ref{tab:constr} reports the comparison two ways. Panel A is a paired test:
every hyperparameter is held fixed and only the construction varies, so the
difference is attributable to construction alone. It is large, highly significant
($t=5.82$), and favours the proposed construction in fifteen of seventeen matched
cells. Panel B sweeps all six construction variants over the same grid and ranks
them by both their best and their average configuration; the proposed
construction places first on both criteria and the conventional pipeline last.
The ordering is therefore not an artefact of a single fortunate parameterisation.

\begin{table}[!htbp]\centering
\caption{Construction. Both constructions swept over an identical grid of 1,286
configurations.}\label{tab:constr}
\begin{tabular}{lcc}
\toprule
\multicolumn{3}{l}{\textbf{Panel A: paired comparison (hyperparameters fixed)}} \\
\midrule\addlinespace[2pt]
matched cells & \multicolumn{2}{c}{17} \\
one book per factor, self-referential & \multicolumn{2}{c}{1.657} \\
blended scores, cross-sectional, equal weight & \multicolumn{2}{c}{0.992} \\
mean difference & \multicolumn{2}{c}{$\mathbf{+0.665}$} \\
$t$-statistic of difference & \multicolumn{2}{c}{\textbf{5.82}} \\
cells favouring proposed construction & \multicolumn{2}{c}{15 of 17} \\
\midrule\addlinespace[2pt]
\multicolumn{3}{l}{\textbf{Panel B: best-of-grid, all six variants}} \\
\midrule\addlinespace[2pt]
construction & best & mean \\
\midrule\addlinespace[2pt]
self / books / risk parity & \textbf{2.571} & \textbf{1.537} \\
self / blend / equal & 2.468 & 1.288 \\
self / books / equal & 2.296 & 1.510 \\
cross / books / equal & 2.200 & 1.152 \\
cross / books / risk parity & 2.113 & 1.129 \\
cross / blend / equal & 1.971 & 0.995 \\
\addlinespace[2pt]
\bottomrule
\end{tabular}
\end{table}

\subsection{Is it alpha?}
\label{sec:alpha}

A Sharpe ratio establishes that a strategy made money; it does not establish that
the return is unexplained by factors already known to be priced. We therefore
regress the weekly strategy return on three controls constructed from the same
panel, an equal-weight market return, a size book formed on trailing dollar
volume, and a twelve-week momentum book, each formed point-in-time
(Table~\ref{tab:alpha}).

Three features stand out. The market loading is $-0.001$ with $t=-0.12$: the book
is market-neutral \emph{empirically}, not merely by construction. Size and
momentum load negatively and significantly, but explain little; the regression
$R^2$ is $0.074$. And the intercept is large and significant, an annualised
alpha of $27.0\%$, accounting for $97\%$ of the strategy's mean return.

Because weekly strategy returns are frequently autocorrelated, which biases
ordinary $t$-statistics upward, we report Newey--West standard errors alongside
the ordinary ones and quote the former. The correction is modest here, first-order
autocorrelation is $+0.063$, moving the alpha $t$-statistic from $5.03$ to
$4.42$; but it is the conservative figure and it still clears the $t>3.0$
threshold of \citet{harvey2016}.

\begin{table}[!htbp]\centering
\caption{Factor-adjusted performance. Weekly strategy returns regressed on an
equal-weight market factor, a size book and a twelve-week momentum book, all
constructed point-in-time from the same panel.}\label{tab:alpha}
\begin{tabular}{lrrr}
\toprule\addlinespace[2pt]
 & coefficient & $t$ (OLS) & $t$ (Newey--West) \\
\midrule\addlinespace[2pt]
$\alpha$ & \textbf{0.00520} & 5.03 & \textbf{4.42} \\
MKT      & $-0.00128$ & $-0.13$ & $-0.12$ \\
SIZE     & $-0.09260$ & $-3.83$ & $-2.69$ \\
MOM      & $-0.06819$ & $-3.39$ & $-2.23$ \\
\midrule\addlinespace[2pt]
annualised $\alpha$ & \multicolumn{3}{r}{$\mathbf{+27.05\%}$} \\
$R^{2}$ & \multicolumn{3}{r}{0.0739} \\
alpha share of mean return & \multicolumn{3}{r}{\textbf{97\%}} \\
observations & \multicolumn{3}{r}{281 weeks} \\
\addlinespace[2pt]
\bottomrule
\end{tabular}

\vspace{4pt}
\begin{minipage}{0.88\textwidth}\footnotesize
\emph{Notes.} Newey--West standard errors use four lags. First-order
autocorrelation of the strategy return is $+0.063$. The market coefficient is
statistically indistinguishable from zero, confirming empirically the
dollar-neutrality imposed by construction.
\end{minipage}
\end{table}

\subsection{Factor contributions}

Table~\ref{tab:loo} reports leave-one-out contributions on a matched 282-week
sample, the full book minus each factor, versus the full book. This is the
quantity of interest for a portfolio member, as distinct from standalone
performance.

The trade-size asymmetry factor ranks \textbf{third of eleven on both measures}:
third by leave-one-out contribution here, and third by standalone performance on
a common sample across all eleven books ($1.921$, $t=3.71$, $n=194$ weeks, the
sample bounded by the three positioning factors). Its contribution is positive in
\emph{both} halves of the sample, indicating structure rather than a
single-period artefact. Separately, \textbf{four factors carry negative
contributions}: the book is marginally better without them. We
report the eleven-factor specification because it was fixed a priori; removing
the underperformers post hoc would be a selection decision made with knowledge of
the outcome, and we decline to make it.

\begin{table}[!htbp]\centering
\caption{Factor contributions. Matched 282-week sample in which every single-factor book is present; full
eleven-factor book returns 2.161 ($t=5.03$).}\label{tab:loo}
\begin{tabular}{lcc}
\toprule
factor & book without it & contribution \\
\midrule\addlinespace[2pt]
Q & 2.033 & $+0.128$ \\
TKU & 2.084 & $+0.078$ \\
\textbf{TSKD} & \textbf{2.115} & $\mathbf{+0.046}$ \\
WRspread & 2.123 & $+0.038$ \\
AVOL & 2.135 & $+0.026$ \\
OFI & 2.144 & $+0.018$ \\
TopChg & 2.150 & $+0.011$ \\
CPVv & 2.170 & $-0.008$ \\
RSJ & 2.179 & $-0.017$ \\
Quad & 2.180 & $-0.019$ \\
CPVm & 2.189 & $-0.028$ \\
\addlinespace[2pt]
\bottomrule
\end{tabular}
\end{table}

\section{Robustness}

\subsection{Drawdown, tails, and the 2022 collapse}
\label{sec:drawdown}

Statistical significance answers whether a result is real. It does not answer
whether the strategy could be held. Table~\ref{tab:dd} reports the equity curve's
properties and, separately, its behaviour through the cryptocurrency decline of
late 2021 and 2022.

Three features are worth drawing out. The maximum peak-to-trough loss is
$7.55\%$, against an annualised return of $27.8\%$, and it was recovered within
sixteen weeks. Skewness is $+1.14$: the return distribution is right-tailed,
which is the opposite of the profile most systematic strategies exhibit and the
opposite of the failure mode associated with them. Excess kurtosis of $6.65$
confirms that the distribution is far from normal, which is why we report these
moments rather than relying on the Sharpe ratio alone.

The second panel tests the neutrality claim where it is most likely to fail. A
regression coefficient is an average over a sample; it does not by itself
establish behaviour in a crisis. Across the 61 weeks from November 2021 the
market fell $82.5\%$ and drew down $84.0\%$. The strategy returned $+9.4\%$ over
the same window with a maximum drawdown of $7.6\%$, and its correlation with the
market across those weeks was $+0.045$. Neutrality is therefore not only an
average property but held through the largest dislocation in the sample.

\begin{table}[!htbp]\centering
\caption{Equity-curve properties and behaviour through the 2021--2022 decline.
Panel A covers the full sample; Panel B the 61 weeks from November 2021.}
\label{tab:dd}
\begin{tabular}{lrr}
\toprule\addlinespace[2pt]
\multicolumn{3}{l}{\textbf{Panel A: full sample (281 weeks)}} \\
\midrule\addlinespace[2pt]
cumulative return & \multicolumn{2}{r}{$+327.6\%$} \\
annualised return & \multicolumn{2}{r}{$+27.76\%$} \\
annualised volatility & \multicolumn{2}{r}{$12.82\%$} \\
Sharpe ratio & \multicolumn{2}{r}{2.166} \\
\textbf{maximum drawdown} & \multicolumn{2}{r}{$\mathbf{-7.55\%}$} \\
\quad weeks to recover & \multicolumn{2}{r}{16} \\
worst week / worst month & \multicolumn{2}{r}{$-6.18\%$ / $-4.66\%$} \\
positive weeks & \multicolumn{2}{r}{$65.5\%$} \\
skewness & \multicolumn{2}{r}{$\mathbf{+1.144}$} \\
excess kurtosis & \multicolumn{2}{r}{$+6.651$} \\
\midrule\addlinespace[2pt]
\multicolumn{3}{l}{\textbf{Panel B: November 2021 to December 2022 (61 weeks)}} \\
\midrule\addlinespace[2pt]
 & market & strategy \\
cumulative return & $-82.5\%$ & $\mathbf{+9.4\%}$ \\
annualised volatility & $92.1\%$ & $12.6\%$ \\
Sharpe ratio & $-1.11$ & $0.67$ \\
maximum drawdown & $-84.0\%$ & $\mathbf{-7.6\%}$ \\
positive weeks & $45.9\%$ & $60.7\%$ \\
correlation with market & \multicolumn{2}{r}{$+0.045$} \\
\addlinespace[2pt]
\bottomrule
\end{tabular}

\vspace{4pt}
\begin{minipage}{0.84\textwidth}\footnotesize
\emph{Notes.} The market is the equal-weight return of the tradeable universe.
The ratio of annualised return to maximum drawdown is $3.7$.
\end{minipage}
\end{table}

\subsection{Walk-forward with rebalance-anchor variation}

Configuration is selected on a trailing 104-week window and evaluated on the
following 26 weeks, rolling forward. Every hyperparameter is re-selected on each
training block, including the quintile fraction: pinning any axis to its
full-sample winner would let the test inherit a choice made with knowledge of the
evaluation period, which is the precise failure this procedure exists to detect.
The candidate set is therefore 36 configurations per fold. Because a weekly
strategy also carries an unmeasured exposure to \emph{which weekday} it
rebalances \citep{hoffstein2019}, the panel is rebuilt from hourly bars on each of
seven anchors and the procedure repeated (Table~\ref{tab:anchor}).

Mean out-of-sample Sharpe is 1.867 and \emph{every one of the seven schedules is
profitable}, each on 182 out-of-sample weeks. The dispersion is itself
informative: a study reporting a single schedule could have reported anything
between 1.074 and 2.575 depending on a choice with no economic content. We report
the mean rather than the maximum. Monday, the schedule actually traded, is second
of the seven rather than first, so the traded configuration is not the flattering
one. Five of seven anchors clear the \citet{harvey2016} hurdle of $t>3.0$;
Friday ($t=2.73$) and Thursday ($t=2.01$) do not, and we report them as failures
rather than excluding them.

The selected quintile fraction varies across folds: $0.10$, $0.20$ and $0.30$
are each chosen in at least one block, and the single best-performing fold
selects $0.10$. No one value drives the result, which is the substantive reason
the parameter can be treated as incidental rather than load-bearing.

\begin{table}[!htbp]\centering
\caption{Walk-forward out-of-sample Sharpe by rebalance anchor. Every
hyperparameter, including the quintile fraction, is selected within each training
block; 36 candidates per fold, 182 out-of-sample weeks per anchor. Harvey et al.
(2016) hurdle is $t>3.0$.}
\label{tab:anchor}
\begin{tabular}{lccccccc|c}
\toprule
anchor & Sun & Mon & Wed & Sat & Tue & Fri & Thu & mean \\
\midrule\addlinespace[2pt]
OOS Sharpe & 2.575 & 2.541 & 1.885 & 1.840 & 1.696 & 1.460 & 1.074 & \textbf{1.867} \\
$t$ & 4.82 & 4.75 & 3.53 & 3.44 & 3.17 & 2.73 & 2.01 & \\
clears $t>3$ & yes & yes & yes & yes & yes & \textbf{no} & \textbf{no} & 5 of 7 \\
\addlinespace[2pt]
\bottomrule
\end{tabular}

\vspace{4pt}
\begin{minipage}{0.84\textwidth}\footnotesize
\emph{Notes.} Walk-forward efficiency, out-of-sample Sharpe as a percentage of
the in-sample best, ranges from 61\% (Thursday) to 113\% (Sunday), above the
50\% threshold on every anchor. An earlier version of this test held the quintile
fraction fixed at its full-sample value and reported a mean of 2.068; releasing
that constraint lowers the figure by 0.20 and costs two anchors their $t>3$
status. The lower number is the correct one.
\end{minipage}
\end{table}

\subsection{Fold geometry}

Table~\ref{tab:geom} varies the train and test window lengths across seven
configurations. Every one produces a profitable out-of-sample series and every
one clears the Harvey threshold, on out-of-sample Sharpe ratios spanning $1.998$
to $2.731$. There is no monotone relationship between training-window length and
performance: the longest window tested is mid-range and the second-shortest is
near the top. We read this as the result being insensitive to fold geometry
rather than as evidence about the optimal window, and the geometry used for the
headline result is mid-range rather than maximal.

\begin{table}[!htbp]\centering
\caption{Fold geometry. Monday anchor; only the train/test split varies. All
hyperparameters, including the quintile fraction, are re-selected within each
training block.}
\label{tab:geom}
\begin{tabular}{cccccc}
\toprule
train & test & folds & $n_{\mathrm{oos}}$ & OOS Sharpe & $t$ \\
\midrule\addlinespace[2pt]
78 & 13 & 16 & 208 & 2.609 & 5.22 \\
78 & 26 & 8 & 208 & 2.525 & 5.05 \\
104 & 13 & 14 & 182 & 2.578 & 4.82 \\
\textbf{104} & \textbf{26} & \textbf{7} & \textbf{182} & \textbf{2.541} & \textbf{4.75} \\
104 & 52 & 3 & 156 & 2.731 & 4.73 \\
130 & 26 & 6 & 156 & 1.998 & 3.46 \\
156 & 26 & 5 & 130 & 2.690 & 4.25 \\
\addlinespace[2pt]
\bottomrule
\end{tabular}
\end{table}

\subsection{Held-out universe}

Time-series testing does not establish that hyperparameters generalise across
\emph{assets}. We hold out symbols: configuration is selected on 78 contracts and
evaluated on 34 that had no influence on the choice, across eight random splits,
with the cross-section rebuilt within each side so ranks, quintile boundaries and
the cost model are computed only among symbols present.

A naive comparison is misleading here, and we flag it because we initially made
it. A 34-contract book is mechanically inferior to a 78-contract one, since
information ratio scales with the square root of breadth \citep{grinold1989}.
Table~\ref{tab:holdout} therefore adds a matched-breadth control: the same
selected configuration evaluated on a random 34 of the \emph{training} contracts.

All eight splits are profitable out of universe. More informative than the ratio
is the stability of the selection: all eight chose identical funding weight,
turnover treatment and quintile width, and five of eight the same ranking window;
fitted noise would not converge this way. We note that only one of eight held-out
splits reaches $t>3.0$, a consequence of reduced breadth, and that the held-out
and control distributions overlap substantially. The defensible claim is that
retention is materially above 50\%, not precisely 83\%.

\begin{table}[!htbp]\centering
\caption{Held-out universe, eight random splits. The matched-breadth control
evaluates the same selected configuration on a random 34 of the \emph{training}
contracts, isolating overfitting from the mechanical breadth effect.}
\label{tab:holdout}
\begin{tabular}{lccc}
\toprule
 & train (78) & held-out (34) & control: seen (34) \\
\midrule\addlinespace[2pt]
mean annualised Sharpe & 2.204 & 1.127 & 1.366 \\
median & 2.169 & 1.194 & 1.407 \\
standard deviation & 0.212 & 0.286 & 0.281 \\
splits profitable & , & \textbf{8 of 8} & , \\
\midrule\addlinespace[2pt]
held-out / train & \multicolumn{3}{l}{51\% \quad \emph{conflates two effects; not interpretable}} \\
control / train & \multicolumn{3}{l}{62\% \quad \emph{breadth, mechanical}} \\
\textbf{held-out / control} & \multicolumn{3}{l}{\textbf{83\%} \quad \textbf{overfitting, at matched breadth}} \\
\addlinespace[2pt]
\bottomrule
\end{tabular}
\end{table}

\subsection{Multiple testing with a measured trial count}

The deflated Sharpe ratio \citep{bailey2014} requires the number of trials and
their dispersion. Both describe the \emph{search} rather than the strategy and
cannot be recovered from the winning backtest, so in practice they are supplied
from recollection, precisely the quantity the correction exists to discipline.
We instead log every evaluated configuration to an append-only registry, written
by the code performing the evaluation, and read both inputs from it
(Table~\ref{tab:dsr}).

We report both readings. The defensible choice is $N=211$, since roughly 1,075 of
the logged cells constitute the construction ablation executed \emph{after} the
strategy was specified, to measure the baseline fairly rather than to select among
candidates. The conservative reading falls below 0.95 and we state it. Its
instability is itself the argument against it: the value fell from 0.9507 to
0.9475 purely because an ablation continued running. A statistic that penalises
conducting additional controlled experiments is not measuring selection bias.

\begin{table}[!htbp]\centering
\caption{Deflated Sharpe ratio. Trial count and dispersion read from an executed
experiment registry rather than supplied by the researcher.}\label{tab:dsr}
\begin{tabular}{lcccc}
\toprule
trial set & $N$ & sd(SR) & $SR^{*}$ & DSR \\
\midrule\addlinespace[2pt]
strategy-search configurations & 211 & 0.0768 & 0.2139 & \textbf{0.9885} \\
all logged configurations & 1{,}286 & 0.0766 & 0.2547 & 0.9475 \\
\addlinespace[2pt]
\bottomrule
\end{tabular}

\vspace{4pt}
\begin{minipage}{0.86\textwidth}\footnotesize
\emph{Notes.} The selected strategy has $T=269$ weekly observations, a weekly
Sharpe of $0.3565$, skewness $+0.145$ and kurtosis $4.391$. Against a zero
benchmark the probabilistic Sharpe ratio is $1.0000$. $SR^{*}$ is the expected
maximum Sharpe across $N$ zero-skill trials.
\end{minipage}
\end{table}

\subsection{Transaction costs}

Table~\ref{tab:cost} scales the entire liquidity-adjusted cost curve by
successive factors of two. The uncapped book is superior only while realised
costs remain close to the modelled level; beyond roughly four times it the
ordering reverses, and by the extreme of the range the uncapped book loses money
while the deployed configuration remains profitable. This is the sense in which
the turnover cap is insurance rather than alpha.

\begin{table}[!htbp]\centering
\caption{Transaction-cost robustness. The entire liquidity-scaled cost curve is
multiplied by the stated factor.}\label{tab:cost}
\begin{tabular}{lcccccc}
\toprule
cost multiple & 1$\times$ & 2$\times$ & 4$\times$ & 8$\times$ & 16$\times$ & 32$\times$ \\
\midrule\addlinespace[2pt]
uncapped & 2.364 & 2.247 & 2.014 & 1.547 & 0.618 & $-1.222$ \\
deployed (capped) & 2.161 & 2.117 & \textbf{2.028} & \textbf{1.851} & 1.496 & 0.789 \\
\addlinespace[2pt]
\bottomrule
\end{tabular}
\end{table}

\subsection{Research--production equivalence}

The relationship between a published backtest and the system trading it is
normally asserted. We test it. Research and live paths call a single shared
implementation, and an automated gate asserts point-in-time integrity,
determinism, dollar-neutrality, non-emptiness and identical position directions,
exiting non-zero so it can block deployment (Table~\ref{tab:gate}).

\begin{table}[!htbp]\centering
\caption{Research--production equivalence gate. Exits non-zero on any failure.}
\label{tab:gate}
\begin{tabular}{lll}
\toprule
assertion & measured & status \\
\midrule\addlinespace[2pt]
point-in-time leak, all factors & $0.00\mathrm{e}{+}00$ & pass \\
determinism & exact & pass \\
dollar-neutrality (max $|$net$|$) & $9.02\mathrm{e}{-}17$ & pass \\
gross exposure bound & $1.462 \le 2.0$ & pass \\
non-emptiness & 308 built, 282 active & pass \\
\textbf{research $=$ live directions} & \textbf{0 / 4{,}480} & \textbf{pass} \\
\addlinespace[2pt]
\bottomrule
\end{tabular}
\end{table}

\section{Economic significance and capacity}

A Sharpe ratio is not an investment. Table~\ref{tab:cap} estimates deployable
capital from realised liquidity, the book's measured breadth (76 concurrent
positions) and measured weekly turnover (0.369 of gross), under a participation
constraint on the 25th-percentile contract by average daily dollar volume, the
binding constraint, since the book must trade its illiquid members as well as its
liquid ones.

The implied figures run from tens to hundreds of millions of dollars depending on
the participation rate assumed. These are order-of-magnitude bounds rather than
precise limits: they assume participation scales linearly with capital and ignore
that market impact is convex in size, so the true ceiling is lower than the linear
extrapolation suggests. They also inherit the maker-fee assumption, which
Section~6.5 shows the deployed configuration survives at eight times the modelled
level, bounding, rather than resolving, the concern.

\begin{table}[!htbp]\centering
\caption{Capacity. Binding constraint is the 25th-percentile contract by average
daily dollar volume (\$3.26m); breadth 76 positions, weekly turnover 0.369.}
\label{tab:cap}
\begin{tabular}{lc}
\toprule
participation rate & implied capacity \\
\midrule\addlinespace[2pt]
1\% & \$33.6m \\
2\% & \textbf{\$67.2m} \\
5\% & \$167.9m \\
10\% & \$335.8m \\
\addlinespace[2pt]
\bottomrule
\end{tabular}

\vspace{4pt}
\begin{minipage}{0.72\textwidth}\footnotesize
\emph{Notes.} Median contract average daily dollar volume is \$6.37m; the tenth
percentile is \$1.93m. The binding constraint is the 25th percentile at \$3.26m.
\end{minipage}
\end{table}

\section{What did not survive}
\label{sec:audit}

\subsection{An audit of our own claims}

Before submission we searched the research codebase for the executable artifact
behind every headline figure. Four had none (Table~\ref{tab:audit}). A separate
claim, that one factor reverses sign between ranking frames, was tested
directly and failed: zero of eleven reverse.

We report this because the corrected figures appear in this paper and the
discarded ones do not, and because the episode motivates the infrastructure in
Sections~6.4 and~6.6.

\begin{table}[!htbp]\centering
\caption{Claims that did not survive audit. All are excluded from this paper.}
\label{tab:audit}
\begin{tabular}{p{0.40\textwidth}p{0.52\textwidth}}
\toprule
claim as previously recorded & what the audit found \\
\midrule\addlinespace[2pt]
walk-forward Sharpe ``2.68--2.95'' & tail-index statistics from a document concerning an unrelated strategy; not Sharpe ratios \\
timing-luck dispersion ``$\pm0.24$'' & no script, output or artifact \\
deflated Sharpe ``0.9966 / 0.9816'' & no artifact \\
conventional baseline ``1.03'' & no artifact; approximately the conventional construction's \emph{mean} (0.992) compared against the proposed construction's \emph{maximum} \\
one factor reverses sign between frames & tested directly; \textbf{zero of eleven} reverse \\
\addlinespace[2pt]
\bottomrule
\end{tabular}
\end{table}

\subsection{Negative results}

We conjectured that crypto factors are structurally more orthogonal than equity
factors, which would make combining them mechanically more valuable here than in
equities. Measured mean pairwise correlation among Alpha101-style books in this
market is $+0.2173$ across 4{,}560 factor pairs. \citet{kakushadze2016}, who
introduced these formulaic alphas, reports an average pairwise correlation of
15.9\% for the same construction in equities. The conjecture is therefore
rejected: factors here are \emph{more} correlated, not less, by roughly a third. The orthogonality that makes per-factor construction
valuable derives from data-source diversity, price, flow, positioning,
trade-size distribution, not from the asset class.

We report Pearson information coefficients alongside Spearman ones throughout.
Returns in this market are severely fat-tailed, and the two can diverge: a factor
may order the typical asset correctly while being wrong about the small number of
observations that dominate realised profit. Rank correlation measures the former,
which is not what a portfolio is paid on. We also report decile monotonicity,
since a genuine cross-sectional sort has a gradient and a positive average
information coefficient alone does not establish one. We report Pearson alongside Spearman throughout.

\section{Conclusion}

Factor investing imported into cryptocurrency perpetual futures inherits an
assumption this market does not support: that its constituents are comparable. We
find the assumption costly but not fatal. Cross-sectional ranking remains
profitable, while ranking each contract against its own history performs
materially better in nine of eleven cases, eight of eleven with carry removed: 
for reasons we partly decompose. We further find that how signals are combined
after ranking matters more than any individual signal, and we exhibit a factor
whose construction depends on information this market publishes exactly and
equity research must infer.

We emphasise the structure of the evidence as much as any single result:
out-of-sample in time and across assets, replicated over seven rebalance schedules
and seven fold geometries, corrected for multiple testing using a trial count read
from an executed log, stress-tested to eight times modelled costs, and verified to
be computed identically by the deployed system. Each test was capable of
overturning the result.

\textbf{Limitations.} Our evidence covers one venue over roughly seven years.
Out-of-sample testing is conducted in time and across the cross-section, but not
across venues. Capacity is estimated under a linear participation assumption that
understates convex impact. One factor's sign convention was inverted with
knowledge of the data. The held-out universe result rests on eight splits. Four of the eleven factors in
our sample contribute negatively on leave-one-out; we retain them because the
specification was fixed before measurement, but a practitioner building a book
rather than testing a method would drop them.

\section*{Data and code availability}

All results in this paper are reproducible from publicly available data. Price,
volume and funding series are drawn from a public exchange bulk archive;
positioning series from a keyless public metrics endpoint. No proprietary or paid
data is used. The analysis code, including the backtest, the walk-forward and
held-out-universe procedures, the experiment registry from which the
multiple-testing correction reads its inputs, and the research--production
equivalence gate, is\ifanon{} available from the author on request and will be
released publicly on acceptance.\else{} available at
\url{https://github.com/dhanyamd/mlops/tree/main/quant_signal}. The replication
set is collected under \texttt{research/code/}, whose README maps each script to
the table in this paper that it produces.\fi

\newpage
\section*{Appendix A: Summary of statistical evidence}

\begin{table}[!htbp]\centering
\caption{Summary of statistical evidence. Every test reported in the paper, the
threshold it is judged against, and the outcome. Three tests do not clear their
threshold and are marked accordingly.}\label{tab:evidence}
\small
\begin{tabular}{p{0.31\textwidth}p{0.24\textwidth}p{0.18\textwidth}c}
\toprule
test & statistic & threshold & outcome \\
\midrule\addlinespace[2pt]
\multicolumn{4}{l}{\textbf{In-sample significance}} \\
$t$-statistic, full book & $t=5.03$, $n=282$ & $t>3.0$ (Harvey) & pass \\
Probabilistic Sharpe vs zero & $\mathrm{PSR}=1.0000$ & $>0.95$ & pass \\
\midrule\addlinespace[2pt]
\multicolumn{4}{l}{\textbf{Multiple-testing correction}} \\
Deflated Sharpe, strategy search & $0.9885$, $N=211$ & $>0.95$ (Bailey) & pass \\
Deflated Sharpe, all logged cells & $0.9475$, $N=1{,}286$ & $>0.95$ & \textbf{fail} \\
\midrule\addlinespace[2pt]
\multicolumn{4}{l}{\textbf{Out-of-sample in time}} \\
Walk-forward, anchor-averaged & $SR=1.867$ & $>0$ & pass \\
Walk-forward, per anchor & $t=2.01$--$4.82$ & $t>3.0$ & 5 of 7 \\
Walk-Forward Efficiency & $103\%$ & $>50$--$60\%$ & pass \\
Fold-geometry robustness & $t=3.46$--$5.22$ & $t>3.0$ & 7 of 7 \\
\midrule\addlinespace[2pt]
\multicolumn{4}{l}{\textbf{Out-of-sample across assets}} \\
Held-out universe, 8 splits & 8 of 8 profitable & $>0$ & pass \\
Retention at matched breadth & $83\%$ & $>50\%$ & pass \\
Held-out per-split $t$ & 1 of 8 reach $t>3.0$ & $t>3.0$ & \textbf{1 of 8} \\
\midrule\addlinespace[2pt]
\multicolumn{4}{l}{\textbf{Rebalance-schedule luck (Hoffstein)}} \\
Timing luck, 7 anchors in-sample & $t=3.67$--$5.03$ & $t>3.0$ & 7 of 7 \\
Anchors profitable & 7 of 7 & , & pass \\
\midrule\addlinespace[2pt]
\multicolumn{4}{l}{\textbf{Construction effect}} \\
Paired difference, 17 cells & $+0.665$, $t=5.82$ & $t>3.0$ & pass \\
Cells favouring construction & 15 of 17 & $>50\%$ & pass \\
\midrule\addlinespace[2pt]
\multicolumn{4}{l}{\textbf{Individual components}} \\
Trade-size asymmetry, standalone & $t=3.35$ & $t>3.0$ & pass \\
Trade-size asymmetry, subperiods & $+0.030$ / $+0.079$ & both $>0$ & pass \\
Carry, standalone & $t=2.91$ & $t>3.0$ & \textbf{fail} \\
\midrule\addlinespace[2pt]
\multicolumn{4}{l}{\textbf{Cost and capacity}} \\
Cost robustness, deployed book & viable to $8\times$ & $\ge 1\times$ & pass \\
Capacity at 2\% participation & \$67m & , & , \\
\midrule\addlinespace[2pt]
\multicolumn{4}{l}{\textbf{Implementation integrity}} \\
Point-in-time leak, all factors & $0.00\mathrm{e}{+}00$ & $=0$ & pass \\
Dollar-neutrality & $9.02\mathrm{e}{-}17$ & $\approx 0$ & pass \\
Research $=$ live directions & 0 of 4{,}480 & $=0$ & pass \\
Determinism & exact & exact & pass \\
\addlinespace[2pt]
\bottomrule
\end{tabular}
\end{table}

\newpage
\section*{Appendix B: Rejected constructions}

The multiple-testing correction in Section~6.4 reads its trial count from an
executed log. Consistency requires that the constructions abandoned during
development be disclosed rather than omitted; a search branch that failed is
still a search branch. Nineteen are enumerated below. All were built, measured
and rejected before the specification reported in this paper was fixed.

\small
\noindent\textbf{Panel A: signal constructions (12)}
\vspace{3pt}
\begin{longtable}{p{0.20\textwidth}p{0.40\textwidth}p{0.32\textwidth}}
\toprule
construction & what it was & why rejected \\
\midrule\addlinespace[2pt]
\endhead
FAS & funding accrual residualised on the price path & measured IC $-0.001$ (Pearson), $+0.002$ (Spearman); $|t|<0.4$, indistinguishable from zero \\
RCGO & capital-gains overhang adapted to crypto & degraded every configuration tested \\
Salience / ST & salience-theory weighting of return distributions & Spearman IC significant, Pearson not; no decile gradient \\
SAT & salience-adjusted turnover & no incremental contribution \\
IFD & smart-money bar selection made sign-aware & subsumed by simpler order-flow imbalance \\
Kyle asymmetry & price impact per unit buy vs sell volume & unstable; requires per-side impact estimates the sample cannot support \\
session-RPV & price--volume correlation split by global session & sessions do not partition crypto flow as trading hours do equities \\
LOP & limit-order pressure proxy & not constructible from public bar data \\
WCC & weighted crowding composite & collinear with existing positioning factors \\
FSD & flow-shape divergence & superseded by the trade-size asymmetry factor \\
TCN$-$ACN & top-count minus all-count normalisation & equivalent to an existing positioning factor after rescaling \\
forward funding cost & expected rather than realised funding & forecast error exceeded the signal \\
\addlinespace[2pt]
\bottomrule
\end{longtable}
\normalsize

\newpage
\small
\noindent\textbf{Panel B: portfolio constructions (7)}
\vspace{3pt}
\begin{longtable}{p{0.20\textwidth}p{0.40\textwidth}p{0.32\textwidth}}
\toprule\addlinespace[2pt]
construction & what it was & why rejected \\
\midrule\addlinespace[2pt]
\endhead
LSN & latent-sector neutralisation via principal components, with the component count set by the Marchenko--Pastur noise edge & removed signal along with the shared mode; no net gain \\
priority turnover & spend the turnover budget on the largest signal changes first & destroyed the uniform-adjustment property that makes the cap work \\
cost-aware turnover & vary adjustment speed by estimated cost & same failure mode as above \\
clientele conditioning & condition weights on positioning composition & no incremental contribution \\
leg specialisation & different construction for the long and short legs & no incremental contribution \\
MOD-correction transfer & transfer a published mispricing correction & does not survive the change of market \\
orthogonalisation transfer & symmetric orthogonalisation of the factor matrix before ranking & no gain once returns rather than scores are combined \\
\addlinespace[2pt]
\bottomrule
\end{longtable}
\normalsize

\newpage
\section*{Pre-submission checklist}
\begin{itemize}\itemsep2pt
\item Cite and differentiate the Management Science crypto factor-pricing
literature; Section~\ref{sec:frame} is a direct empirical response.
\item Insert the AI-use disclosure with the abstract (SSRN requirement).
\item Data and code availability statement (Management Science requirement).
\item Anonymise for Management Science: remove name, affiliation,
acknowledgements; strip PDF metadata.
\item Appendix: nineteen rejected constructions, enumerated.
\item Appendix: 100-word plain-language abstract for the practitioner version.
\end{itemize}

\newpage
\begin{thebibliography}{99}
\bibitem[Asness et al.(2013)]{asness2013} Asness CS, Moskowitz TJ, Pedersen LH (2013) Value and momentum everywhere. \emph{Journal of Finance} 68(3):929--985.
\bibitem[Bailey and L\'opez de Prado(2014)]{bailey2014} Bailey DH, L\'opez de Prado M (2014) The deflated Sharpe ratio. \emph{Journal of Portfolio Management} 40(5):94--107.
\bibitem[Barndorff-Nielsen et al.(2010)]{bns2010} Barndorff-Nielsen OE, Kinnebrock S, Shephard N (2010) Measuring downside risk: Realised semivariance. In \emph{Volatility and Time Series Econometrics} (Oxford University Press), 117--136.
\bibitem[Chordia et al.(2002)]{crs2002} Chordia T, Roll R, Subrahmanyam A (2002) Order imbalance, liquidity, and market returns. \emph{Journal of Financial Economics} 65(1):111--130.
\bibitem[Cong et al.(2026)]{cong2026} Cong LW, Karolyi GA, Tang K, Zhao W (2026) Crypto value, factor pricing, and market segmentation. \emph{Management Science}, Articles in Advance. \url{https://doi.org/10.1287/mnsc.2024.05875}.
\bibitem[Daniel and Titman(1997)]{danieltitman1997} Daniel K, Titman S (1997) Evidence on the characteristics of cross sectional variation in stock returns. \emph{Journal of Finance} 52(1):1--33.
\bibitem[Fama and French(1993)]{famafrench1993} Fama EF, French KR (1993) Common risk factors in the returns on stocks and bonds. \emph{Journal of Financial Economics} 33(1):3--56.
\bibitem[Grinold(1989)]{grinold1989} Grinold RC (1989) The fundamental law of active management. \emph{Journal of Portfolio Management} 15(3):30--37.
\bibitem[Gu et al.(2020)]{gu2020} Gu S, Kelly B, Xiu D (2020) Empirical asset pricing via machine learning. \emph{Review of Financial Studies} 33(5):2223--2273.
\bibitem[Harvey et al.(2016)]{harvey2016} Harvey CR, Liu Y, Zhu H (2016) \ldots and the cross-section of expected returns. \emph{Review of Financial Studies} 29(1):5--68.
\bibitem[Hoffstein et al.(2019)]{hoffstein2019} Hoffstein C, Sibears J, Faber N (2019) Rebalance timing luck: The difference between hired and fired. \emph{Journal of Index Investing} 10(1):27--36.
\bibitem[Kakushadze(2016)]{kakushadze2016} Kakushadze Z (2016) 101 formulaic alphas. \emph{Wilmott} 2016(84):72--81. \url{https://doi.org/10.1002/wilm.10525}.
\bibitem[Kyle(1985)]{kyle1985} Kyle AS (1985) Continuous auctions and insider trading. \emph{Econometrica} 53(6):1315--1335.
\bibitem[Lee and Ready(1991)]{leeready1991} Lee CMC, Ready MJ (1991) Inferring trade direction from intraday data. \emph{Journal of Finance} 46(2):733--746.
\bibitem[Moskowitz et al.(2012)]{moskowitz2012} Moskowitz TJ, Ooi YH, Pedersen LH (2012) Time series momentum. \emph{Journal of Financial Economics} 104(2):228--250.
\bibitem[Patton and Sheppard(2015)]{pattonsheppard2015} Patton AJ, Sheppard K (2015) Good volatility, bad volatility: Signed jumps and the persistence of volatility. \emph{Review of Economics and Statistics} 97(3):683--697.
\bibitem[Zhou(2008)]{zhou2008} Zhou WX (2008) Universal price impact functions of individual trades in an order-driven market. \emph{Quantitative Finance}. \url{https://arxiv.org/abs/0708.3198}.
\end{thebibliography}

\end{document}
